\section{Independent finite-design acquisition verification}
\label{sec:independentfit126}

The independent acquisition audit operates on the four immutable R121 label
files and all 56 stored actors.  It changes neither the original Bellman
objective nor a deployed coefficient.  Forty-eight numerical problems were
originally fitted; the other eight representations are the exactly matched
ordinary-hinge encodings of paid ReLU fits.  The audit reconstructs the
pre-label feature specification, every value row and required one-sided
original-action derivative row, the coefficient box, rounding rule and stored
serialization.  It configures BVLS rather than TRF within the same SciPy
library.  Four actual VALUE solves execute its distinct active-set iterations;
the other 20 return the feasible unconstrained solution through the shared
initial numerical path.  The four VALUE hinge readouts copy the paid ReLU
solutions.  Its independently formulated marginal
program still uses the HiGHS library.  Thus independence here concerns the
implemented feature, problem-assembly, candidate-generation and exact-checking
paths; it does not assert independent authorship or an independently written
linear-programming solver.

The following familiar convex and weak-duality calculations make the checking
path independent of floating solver status.  They are verification lemmas, not
claims that convex duality is a new neural result.  Let
$B=\prod_{j=1}^p[\underline b_j,\overline b_j]$ be a finite rational box,
$a_i,d_i\in\mathbb Q^p$, and $y_i,\ell_i,u_i\in\mathbb Q$, with
$\ell_i\leq u_i$.  The exact finite-design objectives are
\begin{equation}
 F(\theta)=\sum_i(a_i^{\mathsf T}\theta-y_i)^2,
 \qquad
 R(\theta)=\max_i\{d_i^{\mathsf T}\theta-\ell_i,
                         u_i-d_i^{\mathsf T}\theta\}.
\end{equation}
In particular, $R(\theta)\geq0$.  These objectives are conditional on the
original acquired interval labels.  They are not whole-domain approximation
errors or original-policy regret.

\begin{proposition}[Exact finite-box checking bounds]
\label{prop:finitefit126}
For any rational feasible readout $\widehat\theta\in B$, set
$g=\nabla F(\widehat\theta)$ and
\[
 L_F=\max\!\left\{0,F(\widehat\theta)-g^{\mathsf T}\widehat\theta
       +\sum_{j=1}^p\min\{g_j\underline b_j,g_j\overline b_j\}\right\}.
\]
Then $L_F\leq\min_{\theta\in B}F(\theta)\leq F(\widehat\theta)$.
For arbitrary rational weights $w_i^+,w_i^-\geq0$ satisfying
$\sum_i(w_i^++w_i^-)=1$, define
\[
 c=\sum_i(w_i^+-w_i^-)d_i,
 \quad
 L_R=\max\!\left\{0,-\sum_iw_i^+\ell_i+\sum_iw_i^-u_i
        +\sum_{j=1}^p\min\{c_j\underline b_j,c_j\overline b_j\}\right\}.
\]
Then $L_R\leq\min_{\theta\in B}R(\theta)\leq R(\widehat\theta)$.
The conclusions require no numerical stationarity, floating solver success,
strict convexity or uniqueness of the minimizing readout.
\end{proposition}

\begin{proof}
Convexity gives
$F(\theta)\geq F(\widehat\theta)+g^{\mathsf T}(\theta-\widehat\theta)$.
Minimizing this affine function independently over each coordinate of $B$
gives the stated affine lower bound.  The squared loss is nonnegative, so
replacing the lower bound by its maximum with zero preserves validity.
Evaluating the objective at the feasible readout gives the upper bound.

For the marginal problem, the maximum of the $2m$ displayed affine
expressions is at least any convex combination of them.  Therefore, for every
$\theta\in B$,
\[
 R(\theta)\geq
 \sum_iw_i^+(d_i^{\mathsf T}\theta-\ell_i)
 +\sum_iw_i^-(u_i-d_i^{\mathsf T}\theta)
 =c^{\mathsf T}\theta-\sum_iw_i^+\ell_i+\sum_iw_i^-u_i.
\]
Minimizing the affine expression over $B$ yields its coordinatewise box
minimum.  Since at least one of
$d_i^{\mathsf T}\theta-\ell_i$ and $u_i-d_i^{\mathsf T}\theta$ is
nonnegative for every $i$, $R$ is nonnegative.  The maximum with zero is thus
valid.  Feasibility of $\widehat\theta$ gives the upper bound.  Every operation
uses finite rational arithmetic; approximate multipliers enter only as
candidates for the subsequently checked nonnegative unit-mass weights.
\end{proof}

The audit converts a recorded binary64 multiplier to its exact binary
rational value, takes the nonnegative part of its negative, and divides by the
resulting total mass.  If the mass is zero, it uses the predeclared uniform
weights.  Both constructions satisfy Proposition~\ref{prop:finitefit126}
exactly.  Weak witnesses remain weak and are reported without selecting a
more favorable normalization after inspection.  The exact difference between
the two rounded readout objectives is reported separately from the two
optimality-gap upper bounds.

\begin{proposition}[Rounding and what its bound controls]
\label{prop:roundfit126}
Suppose every box endpoint belongs to the lattice $h\mathbb Z$, and let
$P_B$ denote coordinatewise projection onto $B$.  Round a finite raw readout
to the nearest lattice point, with ties to even, and project the result onto
$B$ to obtain $\widehat\theta$.  Then
$\|\widehat\theta-P_B\theta^{\rm raw}\|_\infty\leq h/2$.
For any two feasible readouts $\theta,\theta+\delta\in B$,
\[
 |R(\theta+\delta)-R(\theta)|
       \leq\max_i\|d_i\|_1\,\|\delta\|_\infty,
\]
and, writing $A$ for the matrix with rows $a_i^{\mathsf T}$,
\[
 |F(\theta+\delta)-F(\theta)|
 \leq2\|A\theta-y\|_2\|A\delta\|_2+\|A\delta\|_2^2.
\]
These are finite-design bounds.  They do not supply a continuous-state or
original-policy certificate.
\end{proposition}

\begin{proof}
For a coordinate inside its box interval, nearest-lattice rounding has error
at most $h/2$, and clipping to lattice-valued endpoints cannot increase the
distance to that coordinate.  For a coordinate outside the interval,
monotonicity of nearest-lattice rounding and lattice-valued endpoints imply
that subsequent clipping gives that endpoint.  Its distance from the
projected raw coordinate is zero.  This proves the first assertion,
including the case where a numerical solver slightly violates a box bound.

Each affine expression defining $R$ changes by either
$d_i^{\mathsf T}\delta$ or its negative.  The difference between two finite
maxima is bounded by the largest absolute change in their entries, which is
at most $\max_i\|d_i\|_1\|\delta\|_\infty$.  Finally,
\[
 F(\theta+\delta)-F(\theta)
 =2(A\theta-y)^{\mathsf T}A\delta+\|A\delta\|_2^2.
\]
Taking absolute values and applying Cauchy--Schwarz gives the second bound.
Neither argument takes a supremum outside the declared design.
\end{proof}

The audit implementation uses $h=2^{-30}$, records the actual exact rounded
objectives, and checks their serialized numerator, denominator, binary64 and
hexadecimal identities.  A different coefficient vector can arise from a
nonunique optimization problem; coefficient equality is therefore recorded,
not imposed as a scientific success criterion.  A successful finite-design
audit does not independently reacquire the original labels, verify fitted
error throughout the state--action cube, or create economic observations for
the new readouts.  These remaining limits stay attached to the response to
R97 B8/M9.
